\documentclass{report}
\usepackage{amsmath,amssymb}
\setlength{\parindent}{0mm}
\setlength{\parskip}{1em}
\begin{document}
\begin{center}
\rule{6in}{1pt} \
{\large Greg von Winckel \\
{\bf Reduced and Full Space Methods in Topology Optimization: Applications in Linear Elasticity and Electrical Conduction}}

Sandia National Laboratories \\ MS 1318 \\ Albuquerque \\ NM 87185
\\
{\tt gvonwin@sandia.gov}\\
Eric C.  Cyr\\
Thomas Gardiner\\
Drew Kouri\\
Denis Ridzal\\
John Shadid\end{center}

Topology optimization is a method which aims to compute a spatial
distribution of materials to optimize a
design figure of merit. A standard such problem appears in the design of
truss structures where the physics
are linear elasticity and the intention is to optimize load compliance.
We also consider the problem of
designing electrical conduits in a two port network with topology
optimization. While this appears to be
a straightforward problem, it presents challenges due to an effective
lack of uniqueness of the solution
arising from extreme differences in conductivity in the structure. A
regularization method to obtain
physically desirable solutions is introduced.

Commonly topology optimization problems are solved using mesh-dependent
approaches such as the Method of
Moving Asymptotes which addresses the problem of the inherently concave
objective functions by making a
sequence of convex approximations. With large scale optimization problems
constrained by complex partial
differential equations, more standard optimization approaches such as SQP
and interior point methods have
been demonstrated to have superior convergence properties[1].

In this work we explore solution of elasticity and conductive network
topology optimization problems in the full
space setting with an interior point method using successively penalized
SQP solves and a reduced space method
using trust region methods with bound constraints. The optimization is
carried out using the Rapid Optimization
Library (ROL), a modular C++ package that is part of the Trilinos suite
of scientific computing tools.
Solution of general inequality constrained problems and implementation of
efficient preconditioners is
an area of active development within ROL.

[1] Susana Rojas-Labanda and Mathias Stolpe,
Benchmarking optimization solvers for structural topology
optimization, Struct. Multidisc. Optim. (2015) 52:527�547


\end{document}
